% Generated draft: requires mathematical and source review.
\documentclass[11pt]{article}
\usepackage[margin=25mm]{geometry}
\usepackage{fontspec}
\setmainfont{Latin Modern Roman}
\usepackage{amsmath,amssymb,mathtools,unicode-math}
\setmathfont{Latin Modern Math}
\usepackage{enumitem,needspace,xurl,hyperref}
\setlength{\parindent}{0pt}
\setlength{\parskip}{5pt}
\setlength{\emergencystretch}{3em}
\newcommand{\N}{\mathbb{N}}
\newcommand{\Z}{\mathbb{Z}}
\newcommand{\Q}{\mathbb{Q}}
\newcommand{\R}{\mathbb{R}}
\newcommand{\C}{\mathbb{C}}
\newcommand{\F}{\mathbb{F}}
\newcommand{\A}{\mathbb{A}}
\newcommand{\PP}{\mathbb{P}}
\newcommand{\E}{\mathbb{E}}
\newcommand{\eps}{\varepsilon}
\newcommand{\dd}{\,\mathrm{d}}
\begin{document}
\section*{UnsolvedMath ID: 30005790}
\textbf{Source problem code:} catalogue entry
\subsection{Definitions and mathematical statement}
For a group \(G\), its profinite completion is
\[
\widehat G=\varprojlim_{\substack{N\trianglelefteq G\\[G:N]<\infty}}G/N,
\]
with its inverse-limit topology. Given a specified class \(\mathcal C\) of groups containing \(G\), profinite rigidity of \(G\) within \(\mathcal C\) means
\[
\forall H\in\mathcal C,\qquad \widehat H\cong\widehat G\ \text{as topological groups}\ \Longrightarrow\ H\cong G.
\]
The supplied record does not specify the comparison class or formulate a precise profinite-rigidity conjecture. Its title concerns profinite rigidity, but its substantive excerpt instead states Question 3.1, the Reachability and Coloring Problem, and ends mid-sentence.

The reachability question can be recorded only as the following conditional formulation. Let \(\Gamma=(V,E)\) denote the intended Coxeter galaxy, once its vertices, edges, and vertical and horizontal moves have been defined. A move scheme is represented abstractly by a relation \(R\subseteq V\times V\) of permitted applications. Does there exist a finite family \(\mathcal M\) of permitted vertical and horizontal move schemes such that, for every \(u,v\in V\) lying in the same connected component of \(\Gamma\), there are \(k\in\Z_{\geq0}\), vertices \(v_0,\ldots,v_k\), and schemes \(R_1,\ldots,R_k\in\mathcal M\) satisfying
\[
v_0=u,\qquad v_k=v,\qquad (v_{j-1},v_j)\in R_j\quad(1\leq j\leq k)?
\]
Each application must be an allowed move in the intended galaxy. The excerpt further asks whether an algorithm can output which kinds of moves are needed along a path. Under the interpretation that this means constructing a labelled path, the requested output would consist of such a finite sequence together with the move kinds used at its steps, for every input pair in the same component. The source does not settle whether a stronger classification of necessary move kinds or termination on disconnected input pairs is intended.

\textbf{Unresolved qualifications.} The supplied excerpt does not define the Coxeter galaxy, its permitted presentations or ranks, its edge relation, or either kind of move. Consequently, the displayed formulation is a template, not a fully specified mathematical problem. The record marks the statement as \texttt{unrecoverable} and supplies a verification note saying that the profinite-rigidity phrase is a section heading followed by a proved theorem for ranks at most three. Neither the precise theorem nor a corresponding general conjecture is supplied. No such theorem or conjecture is reconstructed here.
\subsection{Short English statement}
The excerpt asks whether finitely many kinds of vertical and horizontal moves connect every pair of vertices in each component of the Coxeter galaxy, and whether an algorithm can identify the move kinds along a path. The galaxy and moves are undefined in the supplied excerpt, and its relationship to the profinite-rigidity title remains unresolved.
\subsection{Sources}
\begin{itemize}
\item Supplied local UnsolvedMath registry record, permanent ID \(30005790\), source code \texttt{catalogue entry}, titled \emph{Profinite Rigidity in the Coxeter Galaxy}. The record gives proposed year 2024 and no named proposer.
\item \emph{Mini-Worskhop: Artin Groups meet Triangulated Categories} (2024), Oberwolfach Reports, pp.~203--234, DOI \texttt{10.4171/owr/2024/4}: \url{https://doi.org/10.4171/owr/2024/4}. This is the supplied report citation and the literature-source link labelled ``original OWR report.'' The excerpt identifies Question 3.1 as the Reachability and Coloring Problem.
\item Santos Rego--Schwer, \emph{The galaxy of Coxeter groups}, supplied literature-source link: \url{https://doi.org/10.1016/j.jalgebra.2023.12.006}.
\item The excerpt attributes blow-up moves to Howlett and M\"uhlherr, and attributes their sufficiency for vertical navigation to results of Mihalik, Ratcliffe, and Tschantz. It associates the horizontal question with M\"uhlherr's twist conjecture. These are supplied attributions; no separate bibliographic references or precise hypotheses for these claims are supplied.
\item Literature assessment dated exactly \texttt{2026-08-22}: the record labels the problem ``partially solved'' and states that profinite rigidity for Coxeter diagrams of rank at most three is established while the general question remains open. This is supplied dated curation, not an independently verified assessment. The background explicitly requests independent verification of time-sensitive status and cited claims.
\item The supplied statement-verification note says that the official report uses the profinite-rigidity phrase as a section heading followed by a proved theorem for ranks at most three, rather than an open problem. This note is recorded as supplied evidence; the report has not been independently consulted for this draft.
\end{itemize}
\end{document}
